% BHU PYQ -- Professional Exam Template
\documentclass[11pt,a4paper]{article}

\usepackage[a4paper,top=2.5cm,bottom=2.5cm,left=2.8cm,right=2.8cm,headheight=14pt]{geometry}
\usepackage{amsmath,amssymb}
\usepackage[T1]{fontenc}
\usepackage[utf8]{inputenc}
\usepackage{lmodern}
\usepackage{microtype}
\usepackage{fancyhdr}
\usepackage{enumitem}
\usepackage{hyperref}
\usepackage{lastpage}
\usepackage{parskip}

\hypersetup{colorlinks=true,urlcolor=black,linkcolor=black,
  pdftitle={B.Sc. (Hons.) Botany Sem-VI BOB-603 -- BHU PYQ}}

\pagestyle{fancy}
\fancyhf{}
\renewcommand{\headrulewidth}{0.4pt}
\renewcommand{\footrulewidth}{0.4pt}
\fancyhead[L]{\small   \textit{Banaras Hindu University}}
\fancyhead[R]{\small   \textit{Botany Paper: BOB-603}}
\fancyfoot[C]{\small Page  \thepage\ of \pageref{LastPage}}


\newlist{parts}{enumerate}{2}
\setlist[parts,1]{label=(\alph*),leftmargin=*,itemsep=4pt,topsep=3pt}
\setlist[parts,2]{label=(\roman*),leftmargin=*,itemsep=2pt,topsep=2pt}

\newcommand{\pts}[1]{\hfill   \textit{[#1]}}


\begin{document}


\begin{center}
  {\large\bfseries BANARAS HINDU UNIVERSITY}\\[2pt]
  {
\normalsize B.Sc. (Hons.) Semester VI Examination, 2013-14}\\[6pt]
  \rule{0.8\linewidth}{0.5pt}\\[6pt]
  {\large\bfseries Botany}\\[4pt]
  {\large\bfseries Paper: BOB-603 --- Cytogenetics and Evolutionary Processes}\\[4pt]
  {
\normalsize Time: Three hours \hfill Full Marks: 70}
\end{center}

\rule{\linewidth}{0.4pt}

\smallskip



\noindent   \textit{Note: Answer five questions in all, selecting at least two questions each from Section-A and Section-B. Section-C is compulsory. Question No. 9 carries 10 marks and the remaining questions carry 15 marks each. The figures in the right-hand margin indicate marks.}

\bigskip

\begin{center}
   \textbf{SECTION -- A}
\end{center}

\medskip


\noindent   \textbf{Question 1.} Write an account on the ultra structural organization of chromatin.\pts{15}

\medskip


\noindent   \textbf{Question 2.} Define translocation. Describe the cytogenetics of a translocation heterozygote with the help of suitable diagrams.\pts{15}

\medskip


\noindent   \textbf{Question 3.} Define aneuploidy. How will you distinguish between:\pts{15}

\begin{parts}
  \item double monosomics and nullisomics
  \item primary trisomies and secondary trisomics
\end{parts}

\medskip


\noindent   \textbf{Question 4.} What are transitions and transversions? Name few physical mutagens and explain their mode of action.\pts{15}

\bigskip

\begin{center}
   \textbf{SECTION -- B}
\end{center}

\medskip


\noindent   \textbf{Question 5.} Write short notes on any two of the following:\pts{$7.5  \times 2 = 15$}

\begin{parts}
  \item Base substitution
  \item Kappa Particle
  \item Multiple allelism
\end{parts}

\medskip


\noindent   \textbf{Question 6.} Differentiate between any three pairs of the following:\pts{$5  \times 3 = 15$}

\begin{parts}
  \item Allopolyploid and segmental allopolyploidy
  \item Deletion and duplication
  \item Crossing over and linkage
  \item Centromere and Kinetochore
\end{parts}

\medskip


\noindent   \textbf{Question 7.} Write explanatory note on any two of the following:\pts{$7.5  \times 2 = 15$}

\begin{parts}
  \item Heterosis and its evolutionary significance
  \item Quantitative inheritance
  \item Trisomies in   \textit{Datura stramonium}
\end{parts}

\medskip


\noindent   \textbf{Question 8.} Write short notes on any two of the following:\pts{$7.5  \times 2 = 15$}

\begin{parts}
  \item Organellar genetics
  \item Components of Karyotype
  \item Linkage analysis
\end{parts}

\bigskip

\begin{center}
   \textbf{SECTION -- C}
\end{center}

\medskip


\noindent   \textbf{Question 9.} Answer the following in not more than 50 words each:\pts{$1  \times 10 = 10$}

\begin{parts}
  \item Define telomere
  \item What is constitutive heterochromatin?
  \item Define Karyotyping
  \item What is   \textit{Raphanobrassica}?
  \item Role of Colchicine in plants
  \item What is   \textit{in-situ} hybridization?
  \item Define self sterility in plants?
  \item What is Runner's effect?
  \item What do you mean by genetic load?
  \item Define somatic cell hybrids?
\end{parts}

\vfill
\rule{\linewidth}{0.4pt}

\begin{center}\small
\href{https://bhu-exam.vercel.app/}{Click Here}\\[4pt]\href{https://me-aryan.vercel.app/}{Click Here}\\[4pt]\href{https://quashan-me.vercel.app/}{Click Here}
\end{center}

\end{document}
